\chapter{検証環境と方法}\label{chap:3}

本章では、動作検証およびサイクル数評価に用いたシミュレーション環境と3つの比較方式を定義する。
さらに、評価に選定した入力条件、サイクル数の計測区間、およびテストベンチが実行時に行う各種検査項目の設計について述べる。

\section{実行環境と比較条件}\label{sec:3-1}

すべてのシミュレーションにはVivado XSim 2023.1（64-bit、SW Build 3865809）のビヘイビアシミュレーションを用いた。
コンパイルから実行まではxvlog、xelab、xsimの順に行い、テストベンチ上のクロック周期は10 nsに設定した。
ソート対象は符号なし32ビット整数とし、起動時にはinitial処理によって共有RAMおよびレジスタを初期化した状態から実行を開始する。
なお、本CPUは分岐遅延スロットを持たず、jal命令は戻り先アドレスとしてPC+4を保存する仕様である。

並列化の効果を公正に評価するため、表\ref{tab:3}に示す3つの実行方式を比較対象とした。
主たる比較対象は1コア逐次分割方式（serial\_split）と2コア並列方式（dual）である。
1コア条件でも2コア条件と同一のCPUおよび調停器を使用し、第2コア側の要求入力を0に固定することで、ハードウェア構成の差異を排除した。
また、両方式で入力コピー、分割ソート、マージのアルゴリズム構造を一致させた上で、同期や待機を含む総サイクル数を比較する。
一方、配列全体を分割せずに一度にバブルソートするserial\_fullは、マージ処理を伴わないためアルゴリズム構造が異なる参考条件として位置付けた。

\begin{table}[H]
\centering
\caption{比較する実行方式とCPU構成}
\label{tab:3}
\footnotesize
\setlength{\tabcolsep}{3pt}
\renewcommand{\arraystretch}{1.25}
\begin{tabular}{>{\raggedright\arraybackslash}p{\dimexpr0.184267\linewidth-2\tabcolsep\relax}>{\raggedright\arraybackslash}p{\dimexpr0.494717\linewidth-2\tabcolsep\relax}>{\raggedright\arraybackslash}p{\dimexpr0.321016\linewidth-2\tabcolsep\relax}}\toprule
\textbf{方式} & \textbf{ソートとマージの処理内容} & \textbf{CPU構成}\\\midrule
dual & 2区間を並行してバブルソートし、Core0がマージ & 同一CPUを2コア接続\\
serial\_split & 2区間を順にバブルソートし、同一のマージを実行 & 同一CPUを1コア使用\\
serial\_full & 配列全体を一度にバブルソート（マージなし） & 同一CPUを1コア使用\\
\bottomrule\end{tabular}
\end{table}

\section{入力条件と測定区間}\label{sec:3-2}

入力条件として、要素数0、1、2、3、7、9、10、16、17、32、64の範囲から計18種類の入力ケースを選定し、各ケースを上記3方式で実行する計54条件を評価した。
選定したケースには、空配列、単要素、整列済み、逆順、重複要素、全同値、符号なし32ビットの境界値、および長さごとに固定した乱数配列を含めている。
全54条件の内訳と実行結果は付録\ref{chap:B}に示す。

総サイクル数（cycles）は、リセット解除から実行に参加する全コア（1コア構成ではCore0、2コア構成ではCore0およびCore1）が終端ループへ到達するまでのクロック数と定義した。
この区間には、フラグ初期化、入力コピー、ソート、同期待機、マージ、および結果の書き戻しのすべてが含まれる。
処理区間ごとの内訳を把握するため、READY、LEFT\_DONE、RIGHT\_DONE、およびCOMPLETEの各フラグが初めて1として観測されたサイクル数もあわせて記録した。
なお、シミュレーション上のクロック周期10 nsは論理動作の確認用設定であり、実機での動作周波数を示すものではない。

\section{途中状態と最終配列の検査}\label{sec:3-3}

テストベンチには、最終出力の照合に加えて実行途中の異常を検出する検査機構を組み込んだ。
具体的には、ストア命令が受理されるたびにアドレスの未知値（X）や非整列の有無、許容領域外への書き込みを監視する。
2コア実行時には、READY成立前のCore1による書き込み、各コアの担当区間外への書き込み、両側完了前のマージ開始、および他コアの完了フラグへの不正書き込み（代理通知）を検出対象とした。
さらに、ストールが発生したクロックエッジの前後でPCと全32レジスタの値を比較して状態保持を検証し、終了時には配列領域の64ワードと境界ガードが期待値に一致することを照合する。
これらの検査コードの抜粋は付録\ref{chap:C}に示す。なお、ロード命令のアクセスアドレス全般は本検査の対象外である。

検証の妥当性を保つため、テストベンチ側からCPUのPC、レジスタ、RAM、あるいは完了フラグへ値を直接注入することは行わず、リセット後の初期状態からCPU自身の命令実行によって全処理を完了させた。
波形観測用のテストベンチ（tb\_vivado\_sort.v）では、上記の検査機能に加え、入力コピー完了時および終了時の配列スナップショットを波形信号として記録している。
